\section{Threat Model dan Metrik Evaluasi Keamanan}

\textit{Threat model} dalam evaluasi ABAC berbasis \textit{edge} perlu mencakup ancaman yang muncul akibat penggunaan atribut lokal, keterbatasan konektivitas, serta dinamika sinkronisasi atribut dan keputusan otorisasi. NIST SP 800-162 menjelaskan bahwa keputusan akses dalam ABAC sangat bergantung pada validitas, \textit{freshness}, dan kepercayaan terhadap atribut yang digunakan dalam evaluasi kebijakan \parencite{hu_guide_2014}. NIST juga menekankan bahwa perubahan atribut atau \textit{privilege} dapat memengaruhi keputusan akses secara dinamis tanpa perlu mengubah relasi subjek-objek secara langsung. Dalam konteks \textit{edge-IAM}, \textcite{wang_edge-enabled_2025} menegaskan bahwa keputusan akses menjadi dipertanyakan apabila sistem mengasumsikan nilai atribut selalu mutakhir dan berasal dari sumber yang sepenuhnya terpercaya. Selain itu, \textcite{wang_edge-enabled_2025} juga menjelaskan bahwa \textit{local caching} dan sinkronisasi atribut pada \textit{edge} dapat menimbulkan risiko penggunaan \textit{stale attribute} akibat pembaruan atribut yang tertunda atau hilang. Oleh karena itu, \textit{threat model} penelitian ini mencakup \textit{stale attribute, stale role, revocation delay, PIP unavailable, false allow, false deny,} serta \textit{revocation false allow} sebagai ancaman yang relevan pada implementasi ABAC/XACML terdistribusi di lingkungan \textit{edge} IoT.

\textit{False allow} diposisikan sebagai pelanggaran keamanan karena sistem mengizinkan akses yang seharusnya ditolak. Sebaliknya, \textit{false deny} lebih berkaitan dengan \textit{availability} atau \textit{usability} karena sistem menolak akses yang sebenarnya sah. \textit{Revocation false allow} menjadi kasus khusus yang lebih kritis karena hak akses yang telah dicabut masih dapat digunakan akibat keterlambatan sinkronisasi atau penggunaan atribut \textit{cache} yang sudah usang. Ancaman ini konsisten dengan penjelasan \textcite{wang_edge-enabled_2025} mengenai risiko penggunaan attribut usang dalam LAA pada \textit{edge-IAM}. Dengan demikian, \textit{stale cache hit} tidak hanya menjadi indikator performa \textit{cache}, tetapi juga indikator risiko \textit{freshness} yang dapat memengaruhi kualitas dan reliabilitas keputusan otorisasi.

Ancaman \textit{PIP unavailable} digunakan untuk merepresentasikan kondisi ketika sistem tidak dapat memperoleh atribut terbaru dari sumber atribut. Dalam arsitektur ABAC, PDP membutuhkan atribut dan informasi kebijakan dari berbagai komponen untuk melakukan evaluasi akses. \textcite{tall_framework_2023} menjelaskan bahwa keputusan akses dalam ABAC bergantung pada \textit{trust chain} yang melibatkan atribut subjek, objek, lingkungan, dan mekanisme evaluasi kebijakan. Dalam konteks \textit{edge-IAM}, \textcite{wang_edge-enabled_2025} menunjukkan bahwa LAA pada \textit{edge} bergantung pada sinkronisasi atribut dari \textit{cloud} dan dapat menghadapi risiko \textit{stale attribute} ketika konektivitas atau sinkronisasi terganggu. Oleh karena itu, ketika sumber atribut atau PIP tidak tersedia, sistem \textit{edge} harus memilih antara menggunakan atribut lokal yang berpotensi usang atau menolak akses karena informasi tidak cukup untuk membuat keputusan yang reliabel.

Selain itu, \textit{cache poisoning} dan \textit{time manipulation} terhadap TTL dapat diposisikan sebagai ancaman konseptual atau batasan lanjutan. \textcite{yu_dns-sensor_2025} menunjukkan bahwa \textit{cache poisoning} terjadi ketika isi \textit{cache} tidak lagi konsisten dengan sumber otoratif akibat manipulasi respons \textit{cache}. Sementara itu, \textcite{losty_towards_2025} menunjukkan bahwa kegagalan mengikuti TTL dan manipulasi TTL pada sistem DNS IoT dapat memicu ketidakkonsistenan \textit{query} dan meningkatkan risiko \textit{cache poisoning}. Meskipun ancaman tersebut dibahas pada konteks \textit{DNS caching}, konsep yang sama relevan secara konseptual pada \textit{caching} atribut ABAC karena perubahan TTL dapat memengaruhi \textit{freshness} atribut dan konsistensi keputusan otorisasi.

Evaluasi ABAC berbasis \textit{edge} tidak cukup hanya menggunakan metrik performa seperti \textit{average latency} dan \textit{p95 latency}. \textcite{weerasinghe_adaptive_2023} menunjukkan bahwa evaluasi \textit{caching} IoT perlu mempertimbangkan latensi, kemungkinan \textit{delay, hit rate, retrieval cost}, dan \textit{cache-management cost} secara bersamaan untuk memahami \textit{trade-off} antara performa dan efisiensi pengelolaan \textit{cache}. Oleh karena itu, penelitian ini menggunakan metrik performa berupa \textit{average latency} dan \textit{p95 latency}, metrik \textit{overhead} berupa \textit{PIP calls} dan \textit{cache hits}, metrik \textit{freshness} berupa \textit{stale cache hit rate}, serta metrik kualitas keputusan berupa \textit{false allow rate} dan \textit{false deny rate}. Untuk mengevaluasi dampak keterlambatan pencabutan hak akses, digunakan \textit{revocation false allow rate} sebagai indikator risiko ketika akses tetap diizinkan setelah hak akses dicabut.

Dengan demikian, sistem ABAC berbasis \textit{edge} tidak dapat dinilai aman hanya karena menghasilkan latensi rendah atau \textit{cache hit} tinggi. Sistem yang cepat tetapi menghasilkan \textit{false allow} dalam jumlah tinggi tetap tidak dapat dianggap aman karena keputusan otorisasinya gagal melindungi sumber daya. Oleh karena itu, evaluasi perlu menyeimbangkan performa, \textit{overhead, freshness, revocation handling}, kualitas keputusan, serta beban \textit{resource} seperti \textit{CPU time} dan \textit{peak memory}.